\documentclass[10pt,a4paper]{amsart}
\usepackage{amsmath,amssymb,amsthm,graphicx}
\usepackage{verbatim}
\newenvironment{reference}{\comment}{\endcomment}
%\newenvironment{reference}{}{}
\newenvironment{slogan}{\comment}{\endcomment}
\newenvironment{history}{\comment}{\endcomment}

% For commutative diagrams we use Xy-pic
\usepackage[all]{xy}

% We use 2cell for 2-commutative diagrams.
\xyoption{2cell}
% Theorem environments.
%
\theoremstyle{plain}
\newtheorem{theorem}[subsection]{Theorem}
\newtheorem{proposition}[subsection]{Proposition}
\newtheorem{lemma}[subsection]{Lemma}

\theoremstyle{definition}
\newtheorem{definition}[subsection]{Definition}
\newtheorem{example}[subsection]{Example}
\newtheorem{exercise}[subsection]{Exercise}
\newtheorem{situation}[subsection]{Situation}

\theoremstyle{remark}
\newtheorem{remark}[subsection]{Remark}
\newtheorem{remarks}[subsection]{Remarks}

\numberwithin{equation}{subsection}

% Macros
%
\def\lim{\mathop{\mathrm{lim}}\nolimits}
\def\colim{\mathop{\mathrm{colim}}\nolimits}
\def\Spec{\mathop{\mathrm{Spec}}}
\def\Hom{\mathop{\mathrm{Hom}}\nolimits}
\def\Ext{\mathop{\mathrm{Ext}}\nolimits}
\def\SheafHom{\mathop{\mathcal{H}\!\mathit{om}}\nolimits}
\def\SheafExt{\mathop{\mathcal{E}\!\mathit{xt}}\nolimits}
\def\Sch{\mathit{Sch}}
\def\Mor{\mathop{\mathrm{Mor}}\nolimits}
\def\Ob{\mathop{\mathrm{Ob}}\nolimits}
\def\Sh{\mathop{\mathit{Sh}}\nolimits}
\def\NL{\mathop{N\!L}\nolimits}
\def\CH{\mathop{\mathrm{CH}}\nolimits}
\def\proetale{{pro\text{-}\acute{e}tale}}
\def\etale{{\acute{e}tale}}
\def\QCoh{\mathit{QCoh}}
\def\Ker{\mathop{\mathrm{Ker}}}
\def\Im{\mathop{\mathrm{Im}}}
\def\Coker{\mathop{\mathrm{Coker}}}
\def\Coim{\mathop{\mathrm{Coim}}}

% Boxtimes
%
\DeclareMathSymbol{\boxtimes}{\mathbin}{AMSa}{"02}

%
% Macros for moduli stacks/spaces
%
\def\QCohstack{\mathcal{QC}\!\mathit{oh}}
\def\Cohstack{\mathcal{C}\!\mathit{oh}}
\def\Spacesstack{\mathcal{S}\!\mathit{paces}}
\def\Quotfunctor{\mathrm{Quot}}
\def\Hilbfunctor{\mathrm{Hilb}}
\def\Curvesstack{\mathcal{C}\!\mathit{urves}}
\def\Polarizedstack{\mathcal{P}\!\mathit{olarized}}
\def\Complexesstack{\mathcal{C}\!\mathit{omplexes}}
% \Pic is the operator that assigns to X its picard group, usage \Pic(X)
% \Picardstack_{X/B} denotes the Picard stack of X over B
% \Picardfunctor_{X/B} denotes the Picard functor of X over B
\def\Pic{\mathop{\mathrm{Pic}}\nolimits}
\def\Picardstack{\mathcal{P}\!\mathit{ic}}
\def\Picardfunctor{\mathrm{Pic}}
\def\Deformationcategory{\mathcal{D}\!\mathit{ef}}
\usepackage{cleveref}
\crefname{originalreference}{original reference}{original references}
\Crefname{originalreference}{Original reference}{Original references}
\makeatletter
\newlabel{groupoids-lemma-free-action}{{07S2}{source}{}{sourcepointer-1}{}}
\newlabel{groupoids-lemma-free-action@cref}{{[originalreference][0][]07S2}{[1][0][]source}{}{}{}}
\makeatother
\crefname{originalproposition}{original proposition}{original propositions}
\Crefname{originalproposition}{Original proposition}{Original propositions}
\makeatletter
\newlabel{proposition-finite-flat-equivalence-global}{{07S6}{source}{}{sourcepointer-2}{}}
\newlabel{proposition-finite-flat-equivalence-global@cref}{{[originalproposition][0][]07S6}{[1][0][]source}{}{}{}}
\makeatother
\crefname{originalreference}{original reference}{original references}
\Crefname{originalreference}{Original reference}{Original references}
\makeatletter
\newlabel{properties-lemma-ample-finite-set-in-affine}{{01ZY}{source}{}{sourcepointer-3}{}}
\newlabel{properties-lemma-ample-finite-set-in-affine@cref}{{[originalreference][0][]01ZY}{[1][0][]source}{}{}{}}
\makeatother
\crefname{originalreference}{original reference}{original references}
\Crefname{originalreference}{Original reference}{Original references}
\makeatletter
\newlabel{groupoids-definition-principal-homogeneous-space}{{049A}{source}{}{sourcepointer-4}{}}
\newlabel{groupoids-definition-principal-homogeneous-space@cref}{{[originalreference][0][]049A}{[1][0][]source}{}{}{}}
\makeatother
\title{Quotients of schemes by finite locally free free group actions are schemes and torsors}
\author{The Stacks Project Authors}
\date{stacks-project:spaces-properties.tex:local-snapshot-20261012; GNU FDL 1.2 or later}
\begin{document}
\maketitle
\noindent\textit{Editorial note.} All delivered mathematical text is the author's own, copied verbatim from the pinned Stacks Project chapter file spaces-properties.tex; no statement or proof step was rewritten, reordered or removed.\par\medskip
\noindent\textit{Editorial note.} Span recomputation by label pairing (LANE\_BRIEF 10.2): in this chapter the author's proof environment follows the statement environment as a standalone block; the spec covers the complete environment pair (statement span, proof span) as recorded in the pool candidate.\par\medskip
\noindent\textit{Editorial note.} Occupation scan: combined claim\_key + statement\_anchor + generic-line scan over all registration files (active and withdrawn) found 13 S10-geom published cards for spaces-properties.tex (prop-occupied.json, 2026-10-14); those 13 labels are excluded from this batch by design. For the 77 free labels: 0 occupied. The registry mentions field (7557 files) and focused-seeds (645) both point to the same S10-geom cards (dual-source check), no other occupation.\par\medskip
\ifdefined\hypertarget\hypertarget{sourcepointer-1}{}\fi
\noindent\textit{Source reference.} Original reference 07S2: sources/stacks-chow/tags.txt, source lines 9887--9887. Cross-chapter standard prerequisite cited by the author inside the counted span; it is used as a pointer and its own proof is not part of this item.\par\medskip
\ifdefined\hypertarget\hypertarget{sourcepointer-2}{}\fi
\noindent\textit{Source reference.} Original proposition 07S6: sources/stacks-chow/tags.txt, source lines 9891--9891. Cross-chapter standard prerequisite cited by the author inside the counted span; it is used as a pointer and its own proof is not part of this item.\par\medskip
\ifdefined\hypertarget\hypertarget{sourcepointer-3}{}\fi
\noindent\textit{Source reference.} Original reference 01ZY: sources/stacks-chow/tags.txt, source lines 2613--2613. Cross-chapter standard prerequisite cited by the author inside the counted span; it is used as a pointer and its own proof is not part of this item.\par\medskip
\ifdefined\hypertarget\hypertarget{sourcepointer-4}{}\fi
\noindent\textit{Source reference.} Original reference 049A: sources/stacks-chow/tags.txt, source lines 5573--5573. Cross-chapter standard prerequisite cited by the author inside the counted span; it is used as a pointer and its own proof is not part of this item.\par\medskip
\begin{lemma}
\label{lemma-quotient-scheme}
Let $S$ be a scheme. Let $G \to S$ be a group scheme. Let $X \to S$ be
a morphism of schemes. Let $a : G \times_S X \to X$ be an action. Assume that
\begin{enumerate}
\item $G \to S$ is finite locally free,
\item the action $a$ is free,
\item $X \to S$ is affine, or quasi-affine, or projective, or
quasi-projective, or $X$ is isomorphic to an open subscheme of an
affine scheme, or $X$ is isomorphic to an open subscheme of $\text{Proj}(A)$
for some graded ring $A$, or $G \to S$ is radicial.
\end{enumerate}
Then the fppf quotient sheaf $X/G$ is a scheme and $X \to X/G$
is an fppf $G$-torsor.
\end{lemma}

\begin{proof}
We first show that $X/G$ is a scheme. Since the action is free the morphism
$j = (a, \text{pr}) : G \times_S X \to X \times_S X$
is a monomorphism and hence an equivalence relation, see
Groupoids, Lemma \ref{groupoids-lemma-free-action}. The maps
$s, t : G \times_S X \to X$
are finite locally free as we've assumed that $G \to S$ is finite locally
free. To conclude it now suffices to prove the last assumption of
Proposition \ref{proposition-finite-flat-equivalence-global} holds.
Since the action of $G$ is over $S$ it suffices to prove that
any finite set of points in a fibre of $X \to S$ is contained in an
affine open of $X$. If $X$ is isomorphic to an open subscheme of an
affine scheme or isomorphic to an open subscheme of $\text{Proj}(A)$
for some graded ring $A$ this follows from
Properties, Lemma \ref{properties-lemma-ample-finite-set-in-affine}.
If $X \to S$ is affine, or quasi-affine, or projective, or
quasi-projective, we may replace $S$ by an affine open and we
get back to the case we just dealt with. If $G \to S$ is radicial,
then the orbits of points on $X$ under the action of $G$ are singletons
and the condition trivially holds. Some details omitted.

\medskip\noindent
To see that $X \to X/G$ is an fppf $G$-torsor
(Groupoids, Definition \ref{groupoids-definition-principal-homogeneous-space})
we have to show that $G \times_S X \to X \times_{X/G} X$
is an isomorphism and that $X \to X/G$ fppf locally has sections.
The second part is clear from the fact that $X \to X/G$ is surjective
as a map of fppf sheaves (by construction). The first part follows from
the isomorphism $R = U \times_M U$ in the conclusion of
Proposition \ref{proposition-finite-flat-equivalence-global}
(note that $R = G \times_S X$ in our case).
\end{proof}

\end{document}
